% Build from the repo root: tectonic -X compile report/report.tex --outdir .
% Figures come from scripts/pipeline/make_report_figures.py (outputs/figures/*.pdf), so run that first.
\documentclass[10pt,a4paper]{article}

\usepackage[a4paper,top=1.7cm,bottom=1.9cm,left=1.9cm,right=1.9cm]{geometry}
\usepackage{fontspec}
\usepackage{libertinus}
\usepackage{amsmath}
\usepackage{microtype}
\usepackage{xcolor}
\usepackage{graphicx}
\usepackage{booktabs}
\usepackage{tabularx}
\usepackage{array}
\usepackage[font=small,labelfont={bf,color=accent},labelsep=period,skip=4pt]{caption}
\usepackage{subcaption}
\usepackage[round,sort]{natbib}
\usepackage{titlesec}
\usepackage{enumitem}
\usepackage{fancyhdr}
\usepackage[hidelinks]{hyperref}

\definecolor{accent}{HTML}{2A78D6}
\definecolor{ink}{HTML}{2B2B28}
\definecolor{muted}{HTML}{6B6B66}
\definecolor{rule}{HTML}{C9C8C0}

\graphicspath{{../outputs/figures/}{outputs/figures/}}
\setlength{\parindent}{0pt}
\setlength{\parskip}{4.2pt plus 1pt}
\linespread{1.04}
\color{ink}

\titleformat{\section}{\large\bfseries\color{ink}}{\color{accent}\thesection}{0.6em}{}
\titleformat{\subsection}{\normalsize\bfseries\color{ink}}{\color{accent}\thesubsection}{0.5em}{}
\titlespacing*{\section}{0pt}{9pt plus 2pt}{3pt}
\titlespacing*{\subsection}{0pt}{6pt plus 2pt}{2pt}

\pagestyle{fancy}
\fancyhf{}
\renewcommand{\headrulewidth}{0pt}
\fancyfoot[L]{\footnotesize\color{muted}Multi-resolution semantic abstraction over the TKH hypergraph}
\fancyfoot[R]{\footnotesize\color{muted}\thepage}

\newcommand{\code}[1]{\texttt{\small\detokenize{#1}}}
\newcommand{\dn}[1]{\textcolor{muted}{\small DN\,§#1}}
\newcommand{\para}[1]{\textbf{#1}\hspace{0.4em}}
\newcommand{\ci}[2]{[#1,\,#2]}

\begin{document}

\begin{center}
  {\LARGE\bfseries Multi-resolution semantic abstraction\\[2pt] over an evolving knowledge hypergraph\par}
  \vspace{6pt}
  {\color{muted}Ilia Dzneladze \quad·\quad Intern assessment, Constructor Knowledge Labs \quad·\quad September 2026\par}
\end{center}
\vspace{-2pt}

{\color{accent}\rule{\linewidth}{0.6pt}}
\vspace{-6pt}
\begin{quote}\small
I build a three-level hierarchy of concept clusters (12, 50 and 200 groups) for each of four yearly
snapshots of the TKH export and track it through time, using average linkage on one graph that blends an
arity-weighted clique expansion of the hypergraph with a semantic nearest-neighbour graph. Most of my
effort went into evaluating it without grading it by its own ruler. Levels 1 and 2 are coherent to a
blind reader, a blind rater found none of 48 glosses wrong, and routing through the labels beats chance.
Level 0 isn't shown to be coherent, change between snapshots isn't localised, and the hypergraph carries
about 7\% of the affinity, not the 30\% its weight suggests.
\end{quote}
\vspace{-8pt}
{\color{accent}\rule{\linewidth}{0.6pt}}

I designed the method and made the calls described here. A coding agent (Claude Code) wrote, ran and
debugged most of the code under my direction, and \code{AI_USAGE.md} lists what it did. The reasoning
behind each choice, with every number below in more detail, is in \code{DESIGN_NOTES.md}, cited as
\dn{n}. Every follow-up experiment had its keep-or-discard rule written down before it ran \dn{15}, and
I say below where a result went against me.

\section{The data and its snapshots (T1)}

The export has 5,798 nodes of 12 types and 1,429 hyperedges extracted from 52 papers on machine-learned
interatomic potentials. I cut it by hyperedge year: the snapshot at $t$ keeps every edge dated $t$ or
earlier and the nodes those edges reference. Articles and authors stay in as edge members but are never
clustered, since neither is a concept \dn{1}.

\begin{table}[h]
\centering\small
\caption{The four snapshots. \emph{New} counts nodes added since the previous snapshot; \emph{seen
after $t$} counts concept nodes that an early edge references before the corpus first mentions them.}
\label{tab:snapshots}
\begin{tabular}{@{}lrrrrrrr@{}}
\toprule
snapshot & edges & nodes & concepts & papers & new nodes & arity $>2$ & seen after $t$ \\
\midrule
$\le$2020 & 405   & 1,527 & 1,406 & 16 &       & 73\% & 22 \\
$\le$2022 & 589   & 2,200 & 2,046 & 23 & 673   & 76\% & 35 \\
$\le$2024 & 1,063 & 4,184 & 3,895 & 39 & 1,984 & 80\% & 19 \\
$\le$2026 & 1,429 & 5,798 & 5,428 & 52 & 1,614 & 80\% & 0  \\
\bottomrule
\end{tabular}
\end{table}

The structure really is higher-order: median arity is 3 throughout and the largest edge grows from 45 to
65 concepts. The corpus nearly doubles between 2022 and 2024, which matters for stability later. Three
data problems shape the rest. \code{origin_year} is missing for 86\% of nodes, so nothing can depend on
it. There is no entity resolution, so 307 surface forms repeat among concept nodes (``MAE'' is a
separate node in every paper that uses it). And because membership follows the edge year, an early edge
can pull in a node the corpus hasn't seen yet. I first claimed the labeller's input was temporally honest
by construction, and it isn't (§\ref{sec:faith}).

\section{Formal statement (Deliverable 0)}

\para{Construction.} For a snapshot, let $V$ be its concept nodes, $e_c \subseteq V$ the concept members
of hyperedge $e$, $n_e = |e_c|$, and $x_i$ the unit-norm MPNet embedding of node $i$'s surface form.
For $i \ne j$,
\begin{align}
S_{ij} &= \sum_{e \,:\, i,j \in e_c} \frac{1}{n_e - 1}, \qquad
M_{ij} = \cos(x_i, x_j)\cdot \mathbf{1}\!\left[\, j \in \mathrm{kNN}_{15}(i) \ \lor\ i \in \mathrm{kNN}_{15}(j) \,\right], \\
A_{ij} &= \alpha\, \tilde S_{ij} + (1-\alpha)\, \tilde M_{ij}, \quad \alpha = 0.3, \qquad
d(C, D) = 1 - \frac{1}{|C|\,|D|} \sum_{i \in C,\, j \in D} A_{ij},
\end{align}
where $\tilde S = \min\!\big(S / q_{0.99}(S),\, 1\big)$, likewise $\tilde M$, and an entry missing from
the sparse graph counts as $A_{ij} = 0$. Average linkage repeatedly merges the pair of groups with the
smallest $d$, and level $k$ is the cut of the resulting dendrogram with $N_k$ groups,
$(N_0, N_1, N_2) = (12, 50, 200)$. $S$ is a weighted clique expansion: every member spreads one unit of
weight over its co-members, so an edge of arity $n$ adds $n/2$ in total rather than $\binom{n}{2}$, and one
65-concept table can't swamp the graph. With unit weights, edges above arity 10 would hold 94\% of the
pairwise weight; with $1/(n-1)$ they hold 63\% \dn{3}.

\para{A criterion, not an objective.} Average linkage is greedy and optimises nothing globally, so what
the method satisfies is a criterion: at every step, merge the two groups with the highest mean affinity,
and read level $k$ off that process at $N_k$ groups. What the criterion buys is nesting. With missing
entries read as 0, merge heights never decrease, so every cut refines the one above it and the levels are
laminar by construction. Filling missing entries with anything more optimistic can cause inversions and
break that silently. I checked the implementation instead of trusting it: on random sparse graphs it
matches scipy's dense average linkage in merge heights and cuts, and the same test fails for the weighted
variant (WPGMA), so it isn't passing trivially \dn{8}.

\para{The trade-off, and the alternative I rejected.} One joint graph turns the balance between structure
and meaning into a single number I can sweep, and lets a strong semantic tie override weak structure at
any level. I rejected a two-stage scheme, where structure fixes the coarse level and semantics refines
inside it: two groups on the same topic that happen to share no hyperedge could then never meet at the
top, and in a corpus of 52 papers that's the normal case. I argued this rather than built it, which is a
gap. What I got wrong is what $\alpha$ measures. The normalised graphs sit on different scales (median
entry 0.06 for $\tilde S$, 0.57 for $\tilde M$), so at $\alpha = 0.3$ structure carries \textbf{6.6\%} of
the affinity mass, not 30\% \dn{17}. It matters a little more than that, since 94\% of structural pairs
are ones the $k$-NN graph never proposes, but the honest description of what ships is semantic
nearest-neighbour clustering with the hypergraph as a tie-breaker.

\para{Complexity.} Building $S$ costs $\sum_e n_e^2$ (about 66k pairs at 2026), $M$ is a 15-neighbour
search, and the heap-based merge is roughly $O(E \log E)$ in the stored affinities while neighbour lists
stay sparse. All four snapshots build in a few minutes on a laptop CPU, mostly spent embedding.

\begin{table}[h]
\centering\small
\caption{Which of the brief's properties hold by construction and which only empirically.}
\label{tab:props}
\begin{tabularx}{\linewidth}{@{}l l >{\raggedright\arraybackslash}X@{}}
\toprule
property & status & evidence \\
\midrule
P1 laminar refinement & by construction & monotone merge heights; exact check at every snapshot (\code{validate_hierarchy.py}) \\
P2 size budget & by construction & at most $N_k$ groups per cut; $N_0 = 12$ everywhere \\
P3 semantic coherence & empirical & blind intruder test passes at levels 1 and 2, not at level 0 (§\ref{sec:coh}) \\
P4 hyperedge fidelity & partial & $S$ is a projection; the T4 collapse runs and is written out, but doesn't build the levels (§\ref{sec:time}) \\
P5 temporal stability & empirical, not enforced & cross-snapshot ARI 0.35 / 0.45 / 0.65 by level; change isn't localised (§\ref{sec:stab}) \\
P6 faithful labels & empirical, partial & 0 of 48 glosses rated wrong blind; levels 0 and 1 labelled, level 2 not; input temporally honest for about 98\% of nodes (§\ref{sec:faith}) \\
\bottomrule
\end{tabularx}
\end{table}

\section{Where the method sits}

The weighting comes from the work showing that hypergraph Laplacians reduce to clique and star
expansions \citep{zhou2006,agarwal2006}. I borrow the arity-aware weights and skip the spectral step, so
what I claim is a well-weighted projection, not a hypergraph-native method. For the hierarchy I chose
plain agglomeration over principled hierarchical models. \citet{clauset2008} fit a dendrogram as a
generative model and score it on held-out links, the logic of my held-out hyperedge test. Peixoto's
nested block model \citep{peixoto2014} picks the number of groups per level from the data, which would
remove the fixed cut sizes I suspect cause global churn but gives up the hard budget P2. I rejected
hypergraph modularity \citep{kumar2020} because it doesn't yield a laminar family without repair.
\citet{ruggeri2023} show hypergraph communities are hard to detect when pairs rarely recur across edges;
here only 1.7\% of concept pairs recur across papers, which I think is why structure generalises weakly.
\citet{loukas2019} coarsens pairwise graphs, so never has to decide what an edge spanning three groups
becomes, the exact question T4 answers, and SHyPar \citep{sajadinia2024} argues that local coarsening
heuristics like mine lose global structure, a plausible but untested reading of why my multilevel
variant collapsed. For time, the event vocabulary is from \citet{greene2010}, and my scheme is what
\citet{asgari2023} call No-Smoothing: cluster each snapshot, then match. They found it the least smooth
option in most settings; my warm-start test found the smoother option costs coherence on real data.

\section{Time and hyperedge collapse (T3, T4)}
\label{sec:time}

\para{T3.} Each snapshot is clustered on its own and matched to the next by Jaccard overlap on the nodes
present in both. Persistent ids carry forward, and the three transitions log 1,091 events in
\code{temporal_events.json}: 414 births, 268 grows, 128 splits, 108 stable, 103 merges, 43 deaths and 27
shrinks. Births dominate because the concept set grows fourfold. Matching keeps each snapshot
reproducible on its own, at the cost of never pushing the clustering toward stability. Of the three
thresholds I set by feel, only the match threshold is sensitive: moving it from 0.10 to 0.20 takes deaths
from 12 to 78 and splits from 192 to 84 \dn{10}, so every event keeps its raw Jaccard for filtering.
A static picture of levels 0--2 at all four snapshots, with the level-0 flows between them, is in
\code{outputs/figures/hierarchy_across_snapshots.png}.

\para{T4.} If all of a hyperedge's endpoints land in one super-node, the edge becomes internal and is kept
as a count. If they land in two, it becomes a pairwise edge. If they land in three or more, it stays one
$k$-ary edge rather than being exploded into $\binom{k}{2}$ pairs, which would inflate density and turn
one asserted fact into many. Coarse edges over the same super-nodes merge into one weighted edge with a
relation-type breakdown, which is what the rule loses: which original member sat in which original edge.
At level 0 in 2026 the 1,429 edges collapse to 128 pairwise and 464 $k$-ary coarse edges (up to arity
41), plus 7 internal, few because articles pass through as singletons. Exploding the $k$-ary ones would
give 6,926 pairwise edges, about 15 per assertion. The rule runs at every level and is written to
\code{hierarchy.json}, but the shipped levels are cuts of one dendrogram, so nothing reads it back, which
falls short of the brief. The version that does use it, clustering each coarser level on the collapsed
hypergraph below, lost a third or more of level-0 coherence and put about 40\% of nodes in one cluster.
A size-normalised second attempt failed the same way, so under the rule I had set it didn't ship \dn{15}.

\section{Evaluation}

\subsection{Coherence and the circularity hazard}
\label{sec:coh}

Measuring coherence with the MPNet embeddings that built the clusters would grade the method with its
own ruler, so I measured it three other ways, from weakest to strongest. The first is TF-IDF over the same
text, against a random-labels null. The clusters are far from chance ($z$ = 709, 864 and 1,424 by level),
but this is the weakest check, because TF-IDF reads the same strings: 69\% of MPNet neighbour pairs share
a TF-IDF term, against 7\% of random pairs \dn{5}. A degree- and arity-preserving shuffle of the hypergraph
gets almost the same TF-IDF coherence ($z$ = 0.3, 0.4 and 2.9), which is what you'd expect when structure
carries 6.6\% of the mass.

The second check doesn't read text at all. I hide 20\% of the hyperedges, cluster on the rest, and ask how
often the members of a hidden edge share a cluster, relative to a random partition of the same sizes.
At level 0, moving $\alpha$ from 0 to 0.3 to 0.5 raises that lift from 1.95 to 2.42 to 3.53. Holding out
whole papers instead cuts it to 1.98, 2.11 and 2.54, so the structure term mostly encodes co-occurrence
inside papers it has already seen. It clears the bar I set only at level 0 (paired gain \ci{0.01}{0.25}),
and no $\alpha$ beats 0.3 on both axes (Figure~\ref{fig:tradeoff}).

\begin{figure}[!t]
\centering
\includegraphics[width=0.97\linewidth]{structure_meaning_tradeoff.pdf}
\caption{The structure/meaning trade-off across $\alpha$ on the 2026 snapshot: held-out hyperedge lift
against TF-IDF coherence over its null, holding out random edges (blue) or whole papers (orange), 5 seeds.
The ring marks the shipped $\alpha = 0.3$.}
\label{fig:tradeoff}
\end{figure}

The third check is a blind intruder test in the style of \citet{chang2009}. A fresh LLM rater with no
repo access, and a different model from the labeller, sees five members of a super-node plus one
type-matched node from another level-0 branch, and picks the odd one out. Chance is 1 in 6. On null items
built from random nodes it found the intruder once in 20, so it wasn't reading type or length. On real
items it found it 17 of 23 times at level 2 and 14 of 24 at level 1, both well above chance, but only 4 of
12 times at level 0 (\ci{14\%}{61\%}), which fails the rule I set (Figure~\ref{fig:blind}, left). My
conclusion is that levels 1 and 2 hold together for a blind reader and level 0, which is also the least
stable level, is not shown to. With only 12 level-0 groups there is no bigger sample to be had.

\begin{figure}[!t]
\centering
\includegraphics[width=0.84\linewidth]{blind_eval.pdf}
\caption{Blind judgements with Wilson 95\% intervals. Left: intruder detection by level, against chance
and against null items. Right: ratings of glosses shown with their own cluster's unseen members (blue) or
another cluster's (gray).}
\label{fig:blind}
\end{figure}

\subsection{Stability}
\label{sec:stab}

Rebuilding five times with 10\% of hyperedges removed gives ARI 0.58, 0.77 and 0.91 by level. I don't read
that as robustness: the perturbation only touches the structural graph, so it rewards ignoring the
hypergraph (0.75, 0.92 and 0.98 at $\alpha = 0.05$), and I partly chose $\alpha$ on it \dn{17}.
Cross-snapshot ARI on shared nodes is 0.35, 0.45 and 0.65, with intervals from resampling half the shared
nodes (Figure~\ref{fig:stabrout}a). My first interval resampled with replacement and came out entirely
above its own estimate at level 2, because duplicated nodes add agreeing pairs \dn{20}. Per transition,
the level-1 intervals don't overlap: 0.52, then 0.37 over 2022 to 2024 as the corpus nearly doubled, then
0.47. The transitions really differ, and the mean over three is only a summary.

P5 also asks that change be localised to where the corpus changed. It isn't, at least not where new
hyperedges land: at no level does churn rise with a cluster's share of new edges (Spearman $-0.32$,
$+0.03$, $-0.01$, intervals spanning 0), and 351 fine clusters with under 2\% new edges still churn 0.44.
Churn does follow the share of new \emph{semantic} neighbours at levels 1 and 2 (Spearman 0.42 and 0.18),
but I found that only after the planned test failed, so it's a lead. My guess is that fixed cut sizes
force distant clusters to merge whenever a new region needs room. A warm-started variant raised level-2
cross-snapshot ARI from 0.65 to about 0.9 at a real but noisy coherence cost (level-2 $z$ 7\% to 18\%
lower on average), so it didn't ship, though it came closer than I first reported \dn{10}.

\begin{figure}[!t]
\centering
\begin{subfigure}[t]{0.49\linewidth}
  \includegraphics[width=\linewidth]{stability_by_level.pdf}
  \caption{Perturbation (5 seeds) and cross-snapshot stability (half-sampling interval over shared nodes).}
\end{subfigure}\hfill
\begin{subfigure}[t]{0.49\linewidth}
  \includegraphics[width=\linewidth]{routing_vs_chance.pdf}
  \caption{Share of ground-truth nodes kept against the routed pool's share of candidates, with bootstrap
  intervals over 12 questions.}
\end{subfigure}
\caption{Stability by level (a) and coarse-to-fine routing against a random pool of the same size (b).}
\label{fig:stabrout}
\end{figure}

\subsection{Label faithfulness}
\label{sec:faith}

The 248 labels and glosses cover levels 0 and 1; level 2 has none, which is short of P6's ``every
super-node''. They were written by fresh LLM agents from a seeded sample of 25 members plus the type counts
of all members, and the exact prompts are in \code{labeling_input.json}. An earlier set, written by agents
that could see the question files, was replaced \dn{12}. The NLI check grades each gloss against up to 15
members the labeller never saw, with another cluster's members as the control. Contradiction runs 5\% to
14\% by snapshot (controls 45\% to 63\%), and ``not entailed'' runs 60\% to 79\%. I didn't know which of
those was the over-claim rate, so I had a blind rater mark 72 items as accurate, vague or wrong, real and
control mixed (Figure~\ref{fig:blind}, right). It rated \textbf{0 of 48 real glosses wrong}
(\ci{0\%}{7\%}) and 23 of 24 controls wrong, and it called 32 of the 36 real glosses NLI marks ``not
entailed'' accurate. Under the rule I had set, the not-entailed rate measures a weak premise, not
over-claiming. NLI contradiction agrees with the rater only moderately ($\kappa = 0.46$,
\ci{0.24}{0.68}) and raises false alarms. Grading glosses against their members' source-paper titles,
which neither the labeller nor the clustering saw, points the right way (21\% contradicted against 31\% for
controls), but the intervals overlap, so I don't lean on it. Temporal honesty holds for the labeller's
input for about 98\% of nodes, not by construction: the 2020 label \code{L1_S00048} names EquiformerV2,
which the corpus first saw in 2024.

\subsection{Extrinsic: coarse-to-fine retrieval}

For each question, the drill-down routes through label and gloss text (the top 8 of 12 level-0 groups,
then the top 8 of their level-1 children) and ranks the nodes in that pool; the flat baseline ranks all
3,104 method-like nodes by cosine to the question. Ground-truth method names resolve to nodes by exact or
substring match. My first matcher let element nodes (``N'', ``S'', ``Si'') match names like
``physics-informed'', so a third of the answer key was chemical elements, and the numbers looked plausible
the whole time. Requiring both sides of a substring match to be at least four characters leaves 12
questions with a ground truth \dn{16}. Routing works: the pool keeps 58\% of ground-truth nodes while
covering 24\% of candidates, a lift of 0.33 over a same-size random pool (\ci{0.04}{0.54}), and every
budget clears chance. Routing on member centroids instead of labels doesn't (Figure~\ref{fig:stabrout}b),
so the signal is in the labels. Ranking doesn't improve: recall@20 is 0.121 against flat's 0.119
in-sample, and with settings chosen by leave-one-out the paired difference is $-8.7$ points
(\ci{-25.9}{+0.7}, $p = 0.37$), resting on 4 against 5 retrieved ground-truth nodes. The hierarchy finds
the right region, and the ranking inside it is the weak part.

\section{Verified and assumed}

\begin{table}[h]
\centering\small
\begin{tabularx}{\linewidth}{@{}>{\raggedright\arraybackslash}X >{\raggedright\arraybackslash}X@{}}
\toprule
\textbf{Verified}, each with a check or a number behind it & \textbf{Assumed or unresolved} \\
\midrule
Laminarity at every snapshot, by an exact validator & The two-stage alternative: argued, not built \\
The UPGMA implementation, against scipy & Blind judgements from one LLM rater; no second rater, no human pass \\
Coherence against two nulls, held-out hyperedges and a blind intruder test & Why $\alpha$ near 0.3 works, as opposed to that it does \\
Stability intervals by half-sampling, with the old estimator's bias pinned in a test & The temporal thresholds, set by feel; the match threshold is sensitive \\
The NLI judge, against blind gloss ratings & Whether fixed cut sizes cause the global churn \\
The structural share of the affinity (6.6\%), measured & The 7 substring matches in the answer key, unchecked against the papers \\
The ground-truth matcher, fixed and pinned by a test & Level-2 labels, which don't exist \\
A one-command rerun from a bare Python that reproduces all 45 outputs byte for byte & The Linux install, not run end to end \\
\bottomrule
\end{tabularx}
\end{table}

\para{Reproducing it.} \code{py -3.11 scripts/reproduce_all.py} goes from a fresh clone to every result and
figure, building its own environment first. From a folder with no environment at all it reproduced all 45
output files byte for byte against the published checksums in 72 minutes on a laptop CPU. The same text
can embed differently in the last bits on other hardware, so I publish every model output from my run,
and a different result elsewhere can be traced to the embeddings or to the code \dn{23}. A reproducer can
also swap in their own labels and blind ratings and rerun the whole evaluation on them \dn{24}.

\section{With four more weeks}

On the method, I'd normalise the two graphs so that $\alpha$ means influence and re-sweep it, and cut at a
fixed merge height instead of fixed counts, which I suspect drives the global churn. On the evaluation, a
human rater for the intruder and gloss tasks, a semantic-side perturbation so that stability stresses both
signals, and the two-stage alternative measured rather than argued. On the extrinsic task, mapping the
type-B questions and method claims to nodes so that there are more than 12 questions, and then a better
ranker inside the routed pool, since that's where the hierarchy currently loses what routing gains.

{\small
\setlength{\bibsep}{1.5pt plus 0.5pt}
\begin{thebibliography}{11}
\bibitem[Agarwal et~al.(2006)]{agarwal2006} S. Agarwal, K. Branson, S. Belongie. Higher order learning with graphs. \emph{ICML}, 2006.
\bibitem[Asgari et~al.(2023)]{asgari2023} Y. Asgari, R. Cazabet, P. Borgnat. Mosaic benchmark networks: modular link streams for testing dynamic community detection algorithms. arXiv:2310.02840, 2023.
\bibitem[Chang et~al.(2009)]{chang2009} J. Chang, J. Boyd-Graber, C. Wang, S. Gerrish, D. Blei. Reading tea leaves: how humans interpret topic models. \emph{NIPS}, 2009.
\bibitem[Clauset et~al.(2008)]{clauset2008} A. Clauset, C. Moore, M. E. J. Newman. Hierarchical structure and the prediction of missing links in networks. \emph{Nature} 453, 98--101, 2008.
\bibitem[Greene et~al.(2010)]{greene2010} D. Greene, D. Doyle, P. Cunningham. Tracking the evolution of communities in dynamic social networks. \emph{ASONAM}, 176--183, 2010.
\bibitem[Kumar et~al.(2020)]{kumar2020} T. Kumar, S. Vaidyanathan, H. Ananthapadmanabhan, S. Parthasarathy, B. Ravindran. Hypergraph clustering by iteratively reweighted modularity maximization. \emph{Applied Network Science} 5:52, 2020.
\bibitem[Loukas(2019)]{loukas2019} A. Loukas. Graph reduction with spectral and cut guarantees. \emph{JMLR} 20(116), 2019.
\bibitem[Peixoto(2014)]{peixoto2014} T. P. Peixoto. Hierarchical block structures and high-resolution model selection in large networks. \emph{Physical Review X} 4, 011047, 2014.
\bibitem[Ruggeri et~al.(2023)]{ruggeri2023} N. Ruggeri, A. Lonardi, C. De Bacco. Message-passing on hypergraphs: detectability, phase transitions and higher-order information. arXiv:2312.00708, 2023.
\bibitem[Sajadinia et~al.(2024)]{sajadinia2024} H. Sajadinia, A. Aghdaei, Z. Feng. SHyPar: a spectral coarsening approach to hypergraph partitioning. arXiv:2410.10875, 2024.
\bibitem[Zhou et~al.(2006)]{zhou2006} D. Zhou, J. Huang, B. Schölkopf. Learning with hypergraphs: clustering, classification, and embedding. \emph{NIPS}, 2006.
\end{thebibliography}
}

\end{document}
